% ============================================================
%  Paper 01: Axiomatic Foundations of Scale-Dependent Geometric Field Theory
%  Target: RevTeX 4-2 Compilation (SDGFT Series)
% ============================================================
\documentclass[amsmath,amssymb,aps,prd,twocolumn,superscriptaddress,nofootinbib]{revtex4-2}

\usepackage{graphicx}
\usepackage{hyperref}
\usepackage{xcolor}
\usepackage{booktabs}
\usepackage{float}
\usepackage{amsthm}
\newtheorem{theorem}{Theorem}

% ── SDGFT shared macros ──────────────────────────────────────
\usepackage{sdgft-macros}

% ── Document ─────────────────────────────────────────────────
\begin{document}

\preprint{SDGFT-2026-01}

\title{Axiomatic Foundations of Scale-Dependent Geometric Field Theory}

\author{David A. Besemer}
\email{david.besemer@sdgft.org}
\affiliation{Independent Researcher}

\date{\today}

\begin{abstract}
We present the axiomatic foundations of Scale-Dependent Geometric Field Theory (SDGFT), a zero-parameter framework in which all physical couplings and mixing profiles emerge from two topological constants. These axioms, derived from the unique self-dual regular four-dimensional polytope (the 24-cell), are the Fibonacci-lattice conflict Delta = 5/24 and the elementary lattice tension delta = 1/24. We demonstrate how the golden ratio and the six-cone partition of the 3-sphere follow algebraically, setting the stage for a fully coordinate-independent description of spacetime and particle physics.
\end{abstract}

\maketitle

\section{Introduction}
We initiate a program to derive the fundamental parameters of particle physics and cosmology from pure geometry. By identifying the 24-cell as the base unit of discrete spacetime, we show that no free parameters are needed to describe the scale-dependent properties of physical systems.

\section{Theoretical Formulation}
Here we formulate the core mathematical relations governing the system:
\begin{equation}
  \DD = \frac{5}{24}, \quad \dd = \frac{1}{24}
\end{equation}
\begin{equation}
  \DD + \dd = \frac{1}{4} = \sin^2 30^\circ \implies \thetamax = 30^\circ
\end{equation}
\begin{equation}
  \frac{\DD}{\dd} = 5 \implies \gp = \frac{1 + \sqrt{5}}{2}
\end{equation}
\section{The Topological Axioms}
The framework rests on two quantities: the lattice conflict Delta, reflecting Fibonacci self-similarity, and the tension delta. Their sum exactly corresponds to the opening angle of a 30-degree cone, partitioning the three-sphere into six symmetric sectors.

\section{Conclusion}
We have established the topological foundation of SDGFT, providing a rigid framework that links microscopic gauge groups to macroscopic cosmic perturbations.



\begin{thebibliography}{99}
\bibitem{Goldstein_Paper01} D. A. Besemer, ``Axiomatic Foundations of SDGFT,'' SDGFT-2026-01 (in preparation).
\bibitem{Goldstein_Paper04} D. A. Besemer, ``Effective Spectral Dimension and the Banach Fixed-Point Theorem in SDGFT,'' SDGFT-2026-04.
\bibitem{Goldstein_Code} D. A. Besemer, \texttt{sdgft} Python package, \url{https://github.com/david-besemer/sdgft} (2026).
\end{thebibliography}

\end{document}
